\scanpage{194}

\deleted{\unsure{1}} \uline{Prop} Soit $\varphi$ un endom de $\deleted{A_1=k\{X_0,X_1,dX_0,dX_1\}/d\unsure{X_0+}}$ $\mathcal A_1=k\{X_0,X_1\}[dX_0,dX_1]/(dX_0\unsure{+|-}dX_1)_{\deleted{\unsure{pd}}}$, et $\deleted{\unsure{f_1}}\inserted{\varphi^{\unsure{\prime}}}$ un endom de $S$-algèbre différentielle ($S=k\{T\}$) compatible aux deux augmentations $\alpha$ et $\beta$ (qui en fait sont les deux applications simpliciales $\mathcal A_1\to\mathcal A_0$ !). Soit $n\geq0$ ; je dis qu'il existe un et un seul endom de $S$-Alg. différentielle de $\mathcal A_n$, telle que pour les applications simpliciales (\deleted{\unsure{surj}} croissantes !) $\sigma:\Delta_n\unsure{-|\to}\Delta_1$, l'hom. $\mathcal A_*(\sigma):\mathcal A_1\unsure{-|\to}\mathcal A_n$ satisfasse

\[
\deleted{\unsure{\varphi_*\mathcal A_*}}\qquad\varphi_n\mathcal A_*(\sigma)=\mathcal A_*(\sigma)\varphi_1\ .
\]

\deleted{Si $n$} (Si $n\geq1$, on peut se \deleted{\unsure{supposer}} \inserted{borner} dans cette considération aux $\sigma$ croissantes).

% margin note written sideways
\margin{faux tel quel ; vrai ssi $2B_i=0\quad\forall i\geq1$ i.e. $2\varphi^1(X_0^{(i)})=0\quad\forall i\geq1$ (\deleted{\unsure{bien}} \unsure{vide si|évidt si|Ceci est vrai si} $2\,1_k=0$ ; cela signifie, \deleted{\unsure{que|pour}} \deleted{$\unsure{B_i}$} $2.1_k$ \unsure{rég|inv.}$^{\mathrm r}$) que $B_i=0\quad\forall i\geq1$ i.e. $\varphi$ conserve le degré (et alors $\varphi\unsure{\Leftrightarrow|\Rightarrow}$ un \deleted{\unsure{endom}} de $\mathcal A_1^0$ qui soit comp. aux $\alpha^0,\beta^0$).}

\uline{Cor} \deleted{On a $\varphi_1=\varphi$, et} le système $\varphi_*=(\varphi_n)$ est un endom de la $S$-algèbre différentielle semi-simpliciale $\mathcal A_*$, \uline{caractérisé} par le fait que $\varphi_1=\varphi$. L'ens des endom de $S$-alg. diff $\mathcal A_*$ est ainsi en correspondance bi-univoque (conservant les compositions) avec l'ens des $\varphi$ \inserted{qui satisfont aux conditions corrigées plus haut}
% slanted
(à partir de la prop).

\uline{Dém. du corollaire} Il suffit de montrer que si $\tau:\Delta_{n'}\unsure{-|\to}\Delta_n$, d'où $\mathcal A_*(\tau):\mathcal A_n\unsure{-|\to}\mathcal A_{n'}$

\[
\varphi_{n'}\mathcal A_*(\tau)=\mathcal A_*(\tau)\varphi_n\ .
\]

\inserted{C'est clair si $n=0$ car $\mathcal A_0=S$, et par hyp. si $n=1$.} Si on a $n\geq1$ \unsure{et si on|ou bien} \inserted{alors} regarde les $\sigma:\Delta_n\to\Delta_1$ \deleted{\unsure{qui seront}} (on peut se borner aux surjectifs) \inserted{on verra \unsure{que}} les $\mathcal A_*(\sigma):\mathcal A_1\unsure{-|\to}\mathcal A_n$ \deleted{\unsure{ont|sont}} des images qui \uline{engendrent} l'algèbre $\mathcal A_n$, donc il suffit de prouver que

% archive sheet number: 193
